\chapter{Putting it all together: an end-to-end app}
\label{ch:end-to-end}

\begin{aside}
Eighteen chapters of architecture. Now we synthesise. This
chapter builds an end-to-end app --- a small task manager ---
that exercises every concept from Part II. It is the test of
whether the model holds together. Spoiler: it does.
\end{aside}

\section{Specification}

The app: a terminal task manager that:

\begin{itemize}[leftmargin=*]
  \item Loads tasks from a JSON file on startup.
  \item Lets the user add, complete, delete, and edit tasks.
  \item Shows a list on the left, a detail panel on the right.
  \item Filters by status (all / open / done).
  \item Auto-saves to disk after each change.
  \item Supports undo with Ctrl-Z.
  \item Has a help overlay on \code{?}.
\end{itemize}

Six features. Each maps onto a Part II concept.

\section{The model}

\begin{lstlisting}
struct Task {
    int id;
    std::string title;
    std::string notes;
    bool done = false;
    std::chrono::system_clock::time_point created;
};

enum class Filter { All, Open, Done };

enum class Mode { Browsing, Editing, Modal };

struct Model {
    std::vector<Task> tasks;
    int selected = 0;
    Filter filter = Filter::All;
    Mode mode = Mode::Browsing;
    std::optional<int> editing_id;
    std::string edit_buffer;
    std::vector<Model> undo_stack;     // pattern 4
    std::string error;
    bool dirty = false;                 // unsaved
};
\end{lstlisting}

The model is data: vectors, enums, primitives. No pointers, no
shared state, no hidden caches.

\section{The message type}

\begin{lstlisting}
struct Loaded { std::vector<Task> tasks; };
struct LoadFailed { std::string error; };
struct Saved {};
struct SaveFailed { std::string error; };
struct AddTask {};
struct CompleteSelected {};
struct DeleteSelected {};
struct StartEditing {};
struct EditChar { char c; };
struct EditBackspace {};
struct EditCommit {};
struct EditCancel {};
struct CursorUp {};
struct CursorDown {};
struct ChangeFilter { Filter f; };
struct ToggleHelp {};
struct Undo {};
struct Quit {};
using Msg = std::variant<
    Loaded, LoadFailed, Saved, SaveFailed,
    AddTask, CompleteSelected, DeleteSelected,
    StartEditing, EditChar, EditBackspace, EditCommit, EditCancel,
    CursorUp, CursorDown, ChangeFilter, ToggleHelp, Undo, Quit>;
\end{lstlisting}

Eighteen alternatives. Each captures one user intent or system
event. The visitor in \code{update} will have one branch per
alternative; the compiler enforces exhaustiveness.

\section{init}

\begin{lstlisting}
static std::pair<Model, Cmd<Msg>> init() {
    return {
        Model{},
        Cmd<Msg>::task([]() -> Msg {
            auto r = file::read(\"tasks.json\");
            if (!r) return LoadFailed{r.error().context};
            try {
                return Loaded{parse_tasks(r.value())};
            } catch (const std::exception& e) {
                return LoadFailed{e.what()};
            }
        })
    };
}
\end{lstlisting}

The model starts empty; the initial \code{Cmd} kicks off a
file read. The result becomes a \code{Loaded} or \code{LoadFailed}
message. Errors as values (Chapter~\ref{ch:expected}).

\section{update: a tour of the architecture}

A few representative branches.

\subsection*{Loaded: replacing the model}

\begin{lstlisting}
[&](Loaded l) {
    m.tasks = std::move(l.tasks);
    m.selected = 0;
    m.error = \"\";
    return std::pair{m, Cmd<Msg>{}};
}
\end{lstlisting}

Pure replacement. No effect.

\subsection*{AddTask: undo-recorded mutation}

\begin{lstlisting}
[&](AddTask) {
    m.undo_stack.push_back(m);       // pattern 4
    m.tasks.push_back(Task{
        .id = next_id(m.tasks),
        .title = \"New task\",
        .created = std::chrono::system_clock::now()
    });
    m.selected = m.tasks.size() - 1;
    m.dirty = true;
    return std::pair{m, save_cmd(m)};
}
\end{lstlisting}

Three steps: snapshot for undo, mutate, queue a save. The save
command is a \code{Cmd::task} that returns \code{Saved} or
\code{SaveFailed}.

\subsection*{Undo: pattern 4 in action}

\begin{lstlisting}
[&](Undo) {
    if (m.undo_stack.empty()) return std::pair{m, Cmd<Msg>{}};
    auto prev = std::move(m.undo_stack.back());
    m.undo_stack.pop_back();
    prev.undo_stack = std::move(m.undo_stack);  // preserve history
    prev.dirty = true;
    return std::pair{std::move(prev), save_cmd(prev)};
}
\end{lstlisting}

Pop the stack; restore the previous model; queue a save.

\subsection*{EditChar: mode-aware routing}

\begin{lstlisting}
[&](EditChar e) {
    if (m.mode != Mode::Editing) return std::pair{m, Cmd<Msg>{}};
    m.edit_buffer += e.c;
    return std::pair{m, Cmd<Msg>{}};
}
\end{lstlisting}

Mode gates the message: a typed character is meaningful only in
edit mode. Outside of it, ignored.

\subsection*{ToggleHelp: modal stack}

\begin{lstlisting}
[&](ToggleHelp) {
    m.mode = (m.mode == Mode::Modal) ? Mode::Browsing : Mode::Modal;
    return std::pair{m, Cmd<Msg>{}};
}
\end{lstlisting}

A modal is a mode (sorry). The view consults \code{mode} to
draw the help overlay.

\section{view: composing the patterns}

\begin{lstlisting}
static Element view(const Model& m) {
    auto list = render_task_list(m);
    auto detail = render_detail(m);
    auto status = render_status_bar(m);

    auto main = h(list | width<40>, detail | grow<1>);
    auto root = v(main | grow<1>, status | height<1>);

    if (m.mode == Mode::Modal) {
        return overlay(root, render_help());
    }
    return root;
}
\end{lstlisting}

The structure: master-detail (pattern 1) inside a tabbed layout.
The status bar is a footer. The modal overlays everything when
help is open.

\subsection*{render\_task\_list}

\begin{lstlisting}
Element render_task_list(const Model& m) {
    std::vector<Element> rows;
    int idx = 0;
    for (auto& t : m.tasks) {
        if (!filter_match(t, m.filter)) continue;
        bool selected = (idx == m.selected);
        rows.push_back(h(
            text(t.done ? \"[x] \" : \"[ ] \"),
            text(t.title) | (selected ? Bold : Style{})
        ) | (selected ? Bg<60, 65, 80> : Style{}));
        idx++;
    }
    return box(v(rows)) | border<Single>;
}
\end{lstlisting}

A runtime loop builds the rows; the DSL wraps each. The
selected row is highlighted.

\subsection*{Modal overlay}

\begin{lstlisting}
Element overlay(Element bg, Element fg) {
    return z(
        bg | Dim,
        fg | center<Vertical> | center<Horizontal>
    );
}
\end{lstlisting}

\code{z} is a Z-axis stack: children render at the same
position; the last one wins on conflict. Used for tooltips,
modals, dropdowns.

\section{subscribe: pattern by mode}

\begin{lstlisting}
static Sub<Msg> subscribe(const Model& m) {
    auto keys = Sub<Msg>::on_key([&](Key k) -> std::optional<Msg> {
        if (m.mode == Mode::Editing) return edit_key(k);
        if (m.mode == Mode::Modal) return modal_key(k);
        return browse_key(k);
    });
    return keys;
}
\end{lstlisting}

The subscription is mode-aware: each mode dispatches keys
through its own handler. The model decides where input goes.

\section{What this exercises}

\begin{itemize}[leftmargin=*]
  \item Concepts and the \code{Program} contract.
  \item \code{Expected} for file I/O errors.
  \item \code{Cmd::task} for async load and save.
  \item Designated initialisers for tasks.
  \item Variants and \code{visit} for messages.
  \item The DSL for static layout.
  \item Builders for runtime list rows.
  \item Style merge for selection highlighting.
  \item Focus and mode for input routing.
  \item Undo (pattern 4) for history.
  \item Master-detail (pattern 1) for layout.
  \item Modal (overlay) for help.
\end{itemize}

A 200-line app. Every Part II chapter showed up.

\section{Where to take it next}

The app is small; the path forward is large:

\begin{itemize}[leftmargin=*]
  \item Add a filter palette (pattern 8).
  \item Persist undo to disk (so it survives restart).
  \item Add a tag system (pattern 5: incremental search by
    tag).
  \item Sync to a remote backend (pattern 6: side-effecting
    save, optimistic UI).
\end{itemize}

Each is an evening's work because the architecture absorbs the
features without refactoring. That is the win.

\section{The reading order, in retrospect}

If you read Part II in order, you saw the architecture from the
bottom up:

\begin{enumerate}[leftmargin=*]
  \item The Elm loop, conceptually (Ch.~\ref{ch:elm}).
  \item The compile-time DSL that produces elements
    (Ch.~\ref{ch:dsl}).
  \item Elements, the data the renderer consumes
    (Ch.~\ref{ch:elements}).
  \item Layout, the algorithm that places elements
    (Ch.~\ref{ch:layout}).
  \item Concepts, the contracts everything uses
    (Ch.~\ref{ch:concepts}).
  \item Expected, the error-as-value type (Ch.~\ref{ch:expected}).
  \item The pipeline (Ch.~\ref{ch:pipeline}).
  \item Style, theme, palette (Ch.~\ref{ch:style}).
  \item Builders for runtime trees (Ch.~\ref{ch:builders}).
  \item Text and graphemes (Ch.~\ref{ch:text}).
  \item Platform abstraction (Ch.~\ref{ch:platform}).
  \item Focus (Ch.~\ref{ch:focus}).
  \item Scroll (Ch.~\ref{ch:scroll}).
  \item Error boundaries (Ch.~\ref{ch:errorboundary}).
  \item Reactivity preview (Ch.~\ref{ch:reactivity-preview}).
  \item Runtime internals (Ch.~\ref{ch:runtime}).
  \item Testing (Ch.~\ref{ch:testing}).
  \item Widget patterns (Ch.~\ref{ch:widget-arch}).
  \item Compositional patterns (Ch.~\ref{ch:patterns}).
  \item Performance (Ch.~\ref{ch:perf}).
  \item DSL guarantees (Ch.~\ref{ch:dsl-guarantees}).
  \item C++26 features (Ch.~\ref{ch:cpp26-features}).
  \item Comparison to other libraries (Ch.~\ref{ch:comparison}).
\end{enumerate}

Twenty-three chapters. Each is a piece of the architecture. The
end-to-end app above ties them together.

\section{What is next}

Part III walks the rendering pipeline: how the element tree
becomes a canvas of packed cells, how SIMD makes the diff
sub-microsecond, how ANSI is emitted. Part IV walks the
reactive layer in depth and the widget catalogue.

But before turning the page, build something. Even a small
counter. The architecture only feels right once you have lived
in it.

\begin{recap}
A 200-line task manager exercises every architectural concept
from Part II. The pattern: model is data, messages drive state
transitions, the view is a pure function, effects are described
as values. Every concept earned its keep. Part III walks the
rendering details that make the model fast.
\end{recap}
